%%%%%%%%%%%%%%%%%%%%%%%%%%%%%%%%%%%%%%%%%%%%%%%%%%%%%%%%%%%%%%%%%%%%%%%%%%%%%%%%
% ACF-DOC-039.03
% Flood Risk System
%%%%%%%%%%%%%%%%%%%%%%%%%%%%%%%%%%%%%%%%%%%%%%%%%%%%%%%%%%%%%%%%%%%%%%%%%%%%%%%%

\section{Flood Risk System}

\subsection{Executive Summary}

The ACF Flood Risk System is a global hydrological intelligence module
designed to monitor, predict and assess flood hazards using an integrated
Earth system approach.

Flood prediction requires the combination of atmospheric processes,
surface conditions, hydrology and human exposure.

The system integrates:

\begin{itemize}

\item Numerical Weather Prediction

\item Precipitation analysis

\item Hydrological models

\item Satellite observations

\item River monitoring networks

\item Soil moisture analysis

\item Artificial intelligence

\item Earth System Digital Twin

\end{itemize}


The objective is to provide:

\begin{itemize}

\item Global flood forecasting

\item Flash flood warnings

\item River flood prediction

\item Urban flood risk assessment

\item Emergency decision support

\end{itemize}


%%%%%%%%%%%%%%%%%%%%%%%%%%%%%%%%%%%%%%%%%%%%%%%%%%%%%%%%%%%%%%%%%%%%%%%%%%%%%%%%

\subsection{Introduction}

Floods are among the most destructive natural hazards worldwide.

They result from complex interactions between:

\begin{itemize}

\item Atmospheric precipitation

\item Surface runoff

\item River dynamics

\item Soil properties

\item Land use

\item Ocean conditions

\end{itemize}


A flood event can develop from:

\begin{itemize}

\item Extreme rainfall

\item Tropical cyclones

\item Snowmelt

\item Dam failures

\item Storm surge

\item Rapid urban runoff

\end{itemize}


The ACF approach considers floods as a coupled component of the global
Earth system.

%%%%%%%%%%%%%%%%%%%%%%%%%%%%%%%%%%%%%%%%%%%%%%%%%%%%%%%%%%%%%%%%%%%%%%%%%%%%%%%%

\subsection{Global Flood Intelligence Concept}

The ACF flood prediction chain is:

\[
Atmosphere
\rightarrow
Precipitation
\rightarrow
Surface Processes
\rightarrow
Hydrology
\rightarrow
Flood Risk
\rightarrow
Decision Support
\]


The system continuously updates flood probability using:

\begin{itemize}

\item Weather forecasts

\item Hydrological observations

\item Satellite products

\item River measurements

\item Digital Twin simulations

\end{itemize}


%%%%%%%%%%%%%%%%%%%%%%%%%%%%%%%%%%%%%%%%%%%%%%%%%%%%%%%%%%%%%%%%%%%%%%%%%%%%%%%%

\subsection{Hydrological Earth System Modeling}

The hydrological component represents water movement through the
environment.

The water balance equation is:

\[
P-E-R-\Delta S=0
\]


where:

\begin{itemize}

\item $P$ = precipitation

\item $E$ = evapotranspiration

\item $R$ = runoff

\item $\Delta S$ = change in storage

\end{itemize}


The ACF hydrological model represents:

\begin{itemize}

\item Soil water

\item Rivers

\item Lakes

\item Snow storage

\item Groundwater

\item Surface runoff

\end{itemize}


%%%%%%%%%%%%%%%%%%%%%%%%%%%%%%%%%%%%%%%%%%%%%%%%%%%%%%%%%%%%%%%%%%%%%%%%%%%%%%%%
%%%%%%%%%%%%%%%%%%%%%%%%%%%%%%%%%%%%%%%%%%%%%%%%%%%%%%%%%%%%%%%%%%%%%%%%%%%%%%%%
\subsection{Rainfall--Runoff Modeling}

Rainfall--runoff transformation is a fundamental process in flood
prediction.

The ACF hydrological engine converts atmospheric precipitation into
surface water response.

The general relationship is:

\[
Q=f(P,S,T,L)
\]

where:

\begin{itemize}

\item $Q$ = river discharge

\item $P$ = precipitation

\item $S$ = soil state

\item $T$ = temperature conditions

\item $L$ = land surface characteristics

\end{itemize}


The model considers:

\begin{itemize}

\item Rainfall intensity

\item Duration

\item Spatial distribution

\item Soil saturation

\item Terrain slope

\item Vegetation cover

\end{itemize}


%%%%%%%%%%%%%%%%%%%%%%%%%%%%%%%%%%%%%%%%%%%%%%%%%%%%%%%%%%%%%%%%%%%%%%%%%%%%%%%%

\subsection{Surface Runoff Generation}

Surface runoff occurs when precipitation exceeds infiltration capacity.

The infiltration process is represented as:

\[
P=I+R
\]

where:

\begin{itemize}

\item $I$ = infiltration

\item $R$ = runoff

\end{itemize}


The ACF model analyzes:

\begin{itemize}

\item Soil moisture

\item Soil type

\item Land use

\item Surface temperature

\item Vegetation state

\end{itemize}


%%%%%%%%%%%%%%%%%%%%%%%%%%%%%%%%%%%%%%%%%%%%%%%%%%%%%%%%%%%%%%%%%%%%%%%%%%%%%%%%

\subsection{River Flood Forecasting}

The ACF River Flood System predicts river response using hydrological
and hydraulic models.

Important parameters include:

\begin{itemize}

\item River discharge

\item Water level

\item Basin response time

\item Flow velocity

\item Channel capacity

\end{itemize}


The continuity equation is:

\[
\frac{\partial A}{\partial t}
+
\frac{\partial Q}{\partial x}
=
q
\]


where:

\begin{itemize}

\item $A$ = river cross-section area

\item $Q$ = discharge

\item $q$ = lateral inflow

\end{itemize}


%%%%%%%%%%%%%%%%%%%%%%%%%%%%%%%%%%%%%%%%%%%%%%%%%%%%%%%%%%%%%%%%%%%%%%%%%%%%%%%%

\subsection{Hydrological Basin Modeling}

The ACF system divides the planet into hydrological units.

The model includes:

\begin{itemize}

\item Watersheds

\item River networks

\item Reservoirs

\item Lakes

\item Groundwater systems

\end{itemize}


Each basin receives:

\begin{itemize}

\item Precipitation forecast

\item Snowmelt contribution

\item Soil saturation state

\item Historical hydrology

\end{itemize}


%%%%%%%%%%%%%%%%%%%%%%%%%%%%%%%%%%%%%%%%%%%%%%%%%%%%%%%%%%%%%%%%%%%%%%%%%%%%%%%%

\subsection{Flash Flood Detection}

Flash floods are rapid flooding events caused by intense precipitation
over short periods.

The ACF Flash Flood Engine monitors:

\begin{itemize}

\item Extreme rainfall rates

\item Convective precipitation

\item Soil saturation

\item Terrain characteristics

\item Drainage capacity

\end{itemize}


Flash flood risk increases when:

\[
Risk_{flash}
=
Rainfall
\times
SoilSaturation
\times
Terrain
\]


%%%%%%%%%%%%%%%%%%%%%%%%%%%%%%%%%%%%%%%%%%%%%%%%%%%%%%%%%%%%%%%%%%%%%%%%%%%%%%%%

\subsection{Extreme Precipitation Analysis}

The system analyses precipitation extremes using:

\begin{itemize}

\item Rainfall intensity

\item Accumulated precipitation

\item Return periods

\item Anomaly detection

\item Ensemble prediction

\end{itemize}


Important temporal scales:

\begin{itemize}

\item Hourly rainfall

\item Daily accumulation

\item Multi-day events

\end{itemize}


%%%%%%%%%%%%%%%%%%%%%%%%%%%%%%%%%%%%%%%%%%%%%%%%%%%%%%%%%%%%%%%%%%%%%%%%%%%%%%%%

\subsection{Snowmelt Flood Prediction}

In cold regions, floods can result from rapid snowmelt.

The ACF cryosphere-hydrology coupling evaluates:

\begin{itemize}

\item Snow water equivalent

\item Snow temperature

\item Melt rate

\item Frozen soil conditions

\end{itemize}


The melt process is estimated by:

\[
M=f(T,R,S)
\]

where:

\begin{itemize}

\item $M$ = melt rate

\item $T$ = temperature

\item $R$ = radiation

\item $S$ = snow properties

\end{itemize}


%%%%%%%%%%%%%%%%%%%%%%%%%%%%%%%%%%%%%%%%%%%%%%%%%%%%%%%%%%%%%%%%%%%%%%%%%%%%%%%%
%%%%%%%%%%%%%%%%%%%%%%%%%%%%%%%%%%%%%%%%%%%%%%%%%%%%%%%%%%%%%%%%%%%%%%%%%%%%%%%%
\subsection{Urban Flood Modeling}

Urban flooding requires specific modeling because cities modify the
natural hydrological cycle.

The ACF Urban Flood Module integrates:

\begin{itemize}

\item High-resolution terrain models

\item Urban drainage networks

\item Impervious surface mapping

\item Population distribution

\item Infrastructure exposure

\end{itemize}


The urban runoff process depends on:

\[
Q_u=f(P,C,D,T)
\]

where:

\begin{itemize}

\item $Q_u$ = urban runoff

\item $P$ = precipitation

\item $C$ = surface characteristics

\item $D$ = drainage capacity

\item $T$ = terrain properties

\end{itemize}


The system evaluates:

\begin{itemize}

\item Drainage overflow

\item Road flooding

\item Underground infrastructure risk

\item Building exposure

\end{itemize}


%%%%%%%%%%%%%%%%%%%%%%%%%%%%%%%%%%%%%%%%%%%%%%%%%%%%%%%%%%%%%%%%%%%%%%%%%%%%%%%%

\subsection{Coastal Flood and Storm Surge}

Coastal flooding results from the interaction between atmosphere,
ocean and coastal environments.

The ACF coastal flood module combines:

\begin{itemize}

\item Storm surge prediction

\item Sea level anomalies

\item Wave models

\item Tidal information

\item Coastal topography

\end{itemize}


The total coastal water level is:

\[
H_{total}
=
H_{tide}
+
H_{surge}
+
H_{wave}
+
H_{sea}
\]


where:

\begin{itemize}

\item $H_{tide}$ = astronomical tide

\item $H_{surge}$ = storm surge

\item $H_{wave}$ = wave contribution

\item $H_{sea}$ = sea level anomaly

\end{itemize}


%%%%%%%%%%%%%%%%%%%%%%%%%%%%%%%%%%%%%%%%%%%%%%%%%%%%%%%%%%%%%%%%%%%%%%%%%%%%%%%%

\subsection{Satellite Flood Monitoring}

Satellite observation provides global flood monitoring capability.

The ACF observation system integrates:

\begin{itemize}

\item Synthetic Aperture Radar (SAR)

\item Optical satellite imagery

\item Altimetry missions

\item Vegetation monitoring

\item Soil moisture products

\end{itemize}


Satellite flood detection provides:

\begin{itemize}

\item Flood extent mapping

\item Water surface detection

\item Damage assessment

\item Remote region monitoring

\end{itemize}


%%%%%%%%%%%%%%%%%%%%%%%%%%%%%%%%%%%%%%%%%%%%%%%%%%%%%%%%%%%%%%%%%%%%%%%%%%%%%%%%

\subsection{Artificial Intelligence Flood Prediction Engine}

The ACF AI Flood Engine improves prediction capability by combining
physical models and machine learning.

The hybrid prediction equation is:

\[
FloodPrediction=
HydrologicalModel
+
AI Correction
+
Observation
\]


The AI system performs:

\begin{itemize}

\item Flood probability estimation

\item Rainfall pattern recognition

\item River level prediction

\item Satellite image analysis

\item Impact forecasting

\end{itemize}


%%%%%%%%%%%%%%%%%%%%%%%%%%%%%%%%%%%%%%%%%%%%%%%%%%%%%%%%%%%%%%%%%%%%%%%%%%%%%%%%

\subsection{Deep Learning Flood Models}

The AI architecture includes:

\begin{itemize}

\item Convolutional Neural Networks for satellite imagery

\item Recurrent Neural Networks for time series

\item Transformers for long-range prediction

\item Graph Neural Networks for river networks

\end{itemize}


Input data:

\begin{itemize}

\item Rainfall fields

\item River observations

\item Soil moisture

\item Digital elevation models

\item Satellite images

\end{itemize}


Output:

\begin{itemize}

\item Flood probability maps

\item Expected water levels

\item Forecast confidence

\item Impact estimation

\end{itemize}


%%%%%%%%%%%%%%%%%%%%%%%%%%%%%%%%%%%%%%%%%%%%%%%%%%%%%%%%%%%%%%%%%%%%%%%%%%%%%%%%

\subsection{Integration with Earth System Digital Twin}

The Flood Risk System is coupled with the ACF Earth System Digital Twin.

The Digital Twin provides:

\begin{itemize}

\item Atmospheric forcing

\item Land surface conditions

\item Hydrological state

\item Ocean boundary conditions

\item Climate scenarios

\end{itemize}


The coupling allows:

\begin{itemize}

\item Real-time flood simulation

\item Future scenario analysis

\item Climate impact assessment

\item Risk evolution monitoring

\end{itemize}


%%%%%%%%%%%%%%%%%%%%%%%%%%%%%%%%%%%%%%%%%%%%%%%%%%%%%%%%%%%%%%%%%%%%%%%%%%%%%%%%
%%%%%%%%%%%%%%%%%%%%%%%%%%%%%%%%%%%%%%%%%%%%%%%%%%%%%%%%%%%%%%%%%%%%%%%%%%%%%%%%
\subsection{Flood Risk Mapping}

The ACF Flood Risk Mapping System transforms hydrological predictions
into operational risk information.

The risk calculation combines:

\begin{itemize}

\item Hazard probability

\item Flood intensity

\item Exposure

\item Vulnerability

\end{itemize}


The general flood risk equation is:

\[
Risk = Hazard \times Exposure \times Vulnerability
\]


The system generates:

\begin{itemize}

\item Global flood risk maps

\item Regional hazard maps

\item Urban vulnerability maps

\item Infrastructure risk layers

\end{itemize}


%%%%%%%%%%%%%%%%%%%%%%%%%%%%%%%%%%%%%%%%%%%%%%%%%%%%%%%%%%%%%%%%%%%%%%%%%%%%%%%%

\subsection{Exposure and Vulnerability Analysis}

Flood impact depends not only on water presence but also on exposed
elements.

The ACF system integrates:

\begin{itemize}

\item Population density

\item Buildings

\item Transportation networks

\item Energy infrastructure

\item Agricultural areas

\item Critical facilities

\end{itemize}


The vulnerability model considers:

\begin{itemize}

\item Building characteristics

\item Economic value

\item Adaptation capacity

\item Historical flood impacts

\end{itemize}


%%%%%%%%%%%%%%%%%%%%%%%%%%%%%%%%%%%%%%%%%%%%%%%%%%%%%%%%%%%%%%%%%%%%%%%%%%%%%%%%

\subsection{AWCI Operational Flood Dashboard}

The Atmospheric Weather Center Interface (AWCI) receives flood
information from the ACF prediction system.

The operational dashboard displays:

\begin{itemize}

\item Current flood situation

\item Forecast precipitation

\item River levels

\item Flood probability

\item Risk zones

\item Warning status

\end{itemize}


The operational workflow is:

\[
Monitor
\rightarrow
Predict
\rightarrow
Assess
\rightarrow
Alert
\rightarrow
Respond
\]


%%%%%%%%%%%%%%%%%%%%%%%%%%%%%%%%%%%%%%%%%%%%%%%%%%%%%%%%%%%%%%%%%%%%%%%%%%%%%%%%

\subsection{Flood Alert System}

The ACF Alert Engine generates warnings based on dynamic thresholds.

Alert levels include:

\begin{itemize}

\item Green: Normal conditions

\item Yellow: Potential flood risk

\item Orange: Significant flood threat

\item Red: Extreme flood emergency

\end{itemize}


The alert decision considers:

\begin{itemize}

\item Forecast uncertainty

\item Population exposure

\item Expected impact

\item Warning confidence

\end{itemize}


%%%%%%%%%%%%%%%%%%%%%%%%%%%%%%%%%%%%%%%%%%%%%%%%%%%%%%%%%%%%%%%%%%%%%%%%%%%%%%%%

\subsection{Forecast Verification}

The Flood Risk System is continuously evaluated.

Verification metrics include:

\begin{itemize}

\item Flood detection accuracy

\item River level prediction error

\item Spatial flood extent accuracy

\item Warning performance

\end{itemize}


The system uses:

\[
Error =
Observed
-
Predicted
\]


%%%%%%%%%%%%%%%%%%%%%%%%%%%%%%%%%%%%%%%%%%%%%%%%%%%%%%%%%%%%%%%%%%%%%%%%%%%%%%%%

\subsection{Future Development}

Future versions of the ACF Flood Risk System will include:

\begin{itemize}

\item Global kilometer-scale hydrological simulation

\item AI autonomous flood forecasting agents

\item Real-time city-scale digital twins

\item Climate change flood projections

\item Advanced satellite intelligence

\item Fully coupled atmosphere-ocean-hydrology models

\end{itemize}


%%%%%%%%%%%%%%%%%%%%%%%%%%%%%%%%%%%%%%%%%%%%%%%%%%%%%%%%%%%%%%%%%%%%%%%%%%%%%%%%

\subsection{References}

The scientific foundation of the Flood Risk System includes:

\begin{itemize}

\item Hydrology

\item Hydraulic modeling

\item Remote sensing

\item Numerical Weather Prediction

\item Artificial Intelligence

\item Earth System Modeling

\end{itemize}


%%%%%%%%%%%%%%%%%%%%%%%%%%%%%%%%%%%%%%%%%%%%%%%%%%%%%%%%%%%%%%%%%%%%%%%%%%%%%%%%

